\documentclass[11pt]{amsart}
\usepackage{amssymb,mathtools}
\usepackage[margin=1.15in]{geometry}
\usepackage{longtable,array}
\usepackage[hidelinks]{hyperref}

\newtheorem{theorem}{Theorem}[section]
\newtheorem{proposition}[theorem]{Proposition}
\newtheorem{lemma}[theorem]{Lemma}
\newtheorem{corollary}[theorem]{Corollary}
\theoremstyle{definition}
\newtheorem{definition}[theorem]{Definition}
\theoremstyle{remark}
\newtheorem{remark}[theorem]{Remark}
\newcommand{\Zar}{\operatorname{Zar}}
\newcommand{\Spec}{\operatorname{Spec}}
\newcommand{\Hom}{\operatorname{Hom}}
\newcommand{\Idem}{\operatorname{Idem}}
\newcommand{\OO}{\mathcal O}
\newcommand{\units}[1]{#1^{\times}}
\newcommand{\restr}[2]{#1|_{#2}}
\newcommand{\rad}{\sqrt}

\title{Locally ringed lattices, the doubled-origin line,\\
and the projective line}
\date{October 29, 2026}

\begin{document}
\begin{abstract}
We construct the affine spectrum of a commutative ring as a locally
ringed distributive lattice. Localisation is characterised by its
universal property, and the Zariski lattice by the universal property of
support. Sections on an arbitrary lattice element are compatible families
of elements of principal localisations. We prove finite gluing, an
extension theorem for maps out of the structure sheaf, and the existence
of its invertibility support. We identify the global sections of the
affine spectrum with the original ring and prove its universal property
for local morphisms. We define affineness using the canonical map to the
spectrum of global sections and prove that it is equivalent to being
isomorphic to an affine spectrum, and identify principal affine opens
with localisation spectra. We construct the doubled-origin line, compute
its global functions, prove its non-affineness, and classify its points
by idempotents modulo the coordinate. Gluing with the reciprocal
transition gives the projective line, a qcqs scheme over the base ring
with only constant global functions, non-affine over a nonzero base.
The last section records the formalisation status.
\end{abstract}
\maketitle

\section{Conventions}

All rings are commutative with identity; the zero ring is allowed, and
homomorphisms preserve identity. A distributive lattice is bounded, with
bottom $\bot$, top $\top$, meet $\wedge$, and join $\vee$. Lattice maps
preserve both bounds and both binary operations. Empty joins are bottom.

The arguments are constructive. Existence assertions about representatives,
denominators, exponents and finite presentations express mere existence;
they do not specify choices. We use finite choice, which follows by
induction for finite families, and unique choice. In particular, an object
with a uniquely characterising property may be constructed using a merely
existing presentation. No prime ideals or points of a spectrum are used
in the constructions.

The ringed lattice approach to schemes was developed by Coquand,
Lombardi and Schuster in \cite{cls09}; their earlier paper \cite{cls07}
treats the projective spectrum as a distributive lattice. The constructive
Zariski lattice and its algebraic order description are treated in
Coquand and Lombardi \cite{cl03}. We use the locally ringed lattice
formulation of \cite{zeuner}. The affine construction is closely related
to \cite{zm}; the elementary characterisation of localisation is also
used in \cite{isaev}.

\section{Localisation}

Let $S$ be a multiplicative submonoid of a ring $A$.

\begin{definition}
A homomorphism $f:A\to T$ is a \emph{localisation at $S$} if $f(s)$ is a
unit for every $s\in S$ and, for every homomorphism $g:A\to C$ inverting
$S$, there is a unique homomorphism $h:T\to C$ such that $hf=g$.
\end{definition}

The fraction construction gives such a localisation. Its elements are
classes of pairs $(a,s)$ with $s\in S$, where
\[
 (a,s)\sim(b,t)\quad\Longleftrightarrow\quad
 \exists r\in S\; r(at-bs)=0.
\]
Addition and multiplication are given by the usual fraction formulas.

\begin{proposition}\label{prop:elementary}
A homomorphism $f:A\to T$ is a localisation at $S$ if and only if:
\begin{enumerate}
\item $f(s)$ is a unit for every $s\in S$;
\item for every $x\in T$, there merely exist $a\in A$ and $s\in S$ with
      $f(s)x=f(a)$;
\item if $f(a)=0$, there merely exists $s\in S$ with $sa=0$.
\end{enumerate}
\end{proposition}

\begin{proof}
Suppose the three conditions hold and $g:A\to C$ inverts $S$. First,
$f(a)=f(b)$ implies $g(a)=g(b)$: some element of $S$ annihilates $a-b$,
and its image under $g$ is a unit.

For $x\in T$, call $c\in C$ a value of $x$ if, for every fraction equation
$f(s)x=f(a)$, one has $g(s)c=g(a)$. Two fraction equations give
$f(at)=f(bs)$, hence $g(at)=g(bs)$. Cancellation of units shows that the
value calculated from either equation is the same. Thus $x$ has a unique
value, denoted $h(x)$. The fraction rules
\[
 \frac a s+\frac b t=\frac{at+bs}{st},\qquad
 \frac a s\frac b t=\frac{ab}{st}
\]
show that $h$ is a ring homomorphism, and $hf=g$. Any other extension
must give the same values, proving uniqueness.

The fraction construction satisfies the three conditions, the third by
effectiveness of its defining equivalence relation. For any localisation
$f:A\to T$, the universal properties give mutually inverse maps between
$T$ and the fraction ring, commuting with the maps from $A$. Transporting
the three conditions along this isomorphism proves the converse.
\end{proof}

\begin{corollary}\label{cor:agree}
Two homomorphisms out of a localisation which agree on $A$ are equal.
\end{corollary}

Write $A_g$ for a localisation at the powers of $g\in A$, and
$f_g:A\to A_g$ for its structural map. The particular construction of
$A_g$ will not matter until a family of such rings is chosen.

\section{The Zariski lattice}

For a finite family $b=(b_1,\ldots,b_m)$, membership in the ideal $(b)$
means mere existence of coefficients $c_j$ with $x=\sum_j b_jc_j$.
Define
\[
 x\in\rad{(b)}\quad\Longleftrightarrow\quad
 \exists n\geq0\; x^n\in(b).
\]

\begin{lemma}\label{lem:radical}
The radical of a finitely generated ideal is an ideal. Moreover,
$x^n\in\rad{(b)}$ implies $x\in\rad{(b)}$.
\end{lemma}
\begin{proof}
Closure under multiplication by a scalar follows from
$(cx)^n=c^nx^n$. For addition use the elementary identity
\[
 (x+y)^{p+q}=U x^p+V y^q
\]
for suitable $U,V\in A$. It is proved by double induction on $p,q$.
If either exponent is zero the assertion is immediate. In the successor
case, use the two inductive expressions for the common exponent
$p+(q+1)=(p+1)+q$, multiply the first by $x$ and the second by $y$, and
add. This proves the identity without the binomial theorem.
Finally, $(x^n)^m=x^{nm}$ proves the last assertion.
\end{proof}

For finite families put
\[
 a\preceq b\quad\Longleftrightarrow\quad
 \text{every }a_i\text{ belongs to }\rad{(b)}.
\]
This is a preorder. For transitivity, if $a\preceq b\preceq c$, a power
of each $a_i$ is a linear combination of the $b_j$. By
Lemma~\ref{lem:radical} this combination belongs to $\rad{(c)}$, and so
does $a_i$.

\begin{definition}
The \emph{Zariski lattice} $\Zar(A)$ is the set of finite families modulo
the equivalence $a\preceq b$ and $b\preceq a$. Write $[a]$ for a class
and $D(x)=[(x)]$.
\end{definition}

Join is concatenation of families; meet is the family of pairwise
products. The empty family is bottom and $(1)$ is top. Indeed,
concatenation is a least upper bound for $\preceq$. Pairwise products
give a greatest lower bound: if $x^p\in(a)$ and $x^q\in(b)$, then
$x^{p+q}\in(a)(b)$. These facts also establish the lattice laws;
distributivity follows by distributing pairwise products over
concatenation.

\begin{theorem}[Formal Nullstellensatz]\label{thm:null}
For finite families $a,b$,
\[
 [a]\leq[b]\quad\Longleftrightarrow\quad a\preceq b.
\]
Every element of $\Zar(A)$ is merely a finite join of elements $D(x)$,
and $[a]=\bigvee_iD(a_i)$.
\end{theorem}
\begin{proof}
The order of the quotient lattice is characterised by
$[a]\wedge[b]=[a]$. If $a\preceq b$, the greatest-lower-bound property
identifies $a$ with its family of products with $b$. Conversely, this
identification implies $a\preceq ab\preceq b$. Finite generation follows
from the description of join and the existence of representatives.
\end{proof}

\begin{definition}
A \emph{support} of $A$ in a distributive lattice $M$ is a map $d:A\to M$
such that
\[
 d(0)=\bot,\quad d(1)=\top,\quad d(xy)=d(x)\wedge d(y),\quad
 d(x+y)\leq d(x)\vee d(y).
\]
\end{definition}

\begin{theorem}[Universal support]\label{thm:support}
The map $D:A\to\Zar(A)$ is a support. Every support $d:A\to M$ factors
uniquely through a lattice map $\bar d:\Zar(A)\to M$. Its value is
\[
 \bar d([a])=\bigvee_i d(a_i).
\]
\end{theorem}
\begin{proof}
The support identities for $D$ follow from ideal membership and the
lattice operations. For a general support, induction gives
\[
 d(x)\leq d(x^n),\qquad
 d\Bigl(\sum_j b_jc_j\Bigr)\leq\bigvee_jd(b_j).
\]
Thus $a\preceq b$ implies $\bigvee_i d(a_i)\leq\bigvee_jd(b_j)$.
The proposed formula descends to the quotient. It preserves joins by
concatenation and meets by distributivity and $d(xy)=d(x)\wedge d(y)$;
it also preserves both bounds. Uniqueness follows from finite generation.
\end{proof}

\begin{proposition}[The localisation bridge]\label{prop:bridge}
For $g,h\in A$,
\[
 D(h)\leq D(g)
 \quad\Longleftrightarrow\quad
 \exists n,y\;h^n=gy
 \quad\Longleftrightarrow\quad f_h(g)\in\units{A_h}.
\]
\end{proposition}
\begin{proof}
The first equivalence is Theorem~\ref{thm:null}. An equation $h^n=gy$
makes $f_h(g)$ a unit. Conversely, represent an inverse of $f_h(g)$ by
$y/h^n$. Its inverse equation gives $f_h(gy-h^n)=0$. Kernel detection
gives $h^m(gy-h^n)=0$ for some $m$, hence
$h^{m+n}=g(h^my)$.
\end{proof}

If $D(h)\leq D(g)$, the universal property therefore gives a unique map
\[
 \rho_{h,g}:A_g\longrightarrow A_h,\qquad \rho_{h,g}f_g=f_h.
\]
Corollary~\ref{cor:agree} gives identity and composition:
\[
 \rho_{g,g}=1,\qquad \rho_{k,h}\rho_{h,g}=\rho_{k,g}.
\]
These maps are independent of witnesses of the order relations.

\section{Sections on the whole lattice}

Choose a localisation $A_g$ for every $g\in A$. For $u\in\Zar(A)$ define
\[
 \OO_A(u)=\left\{(s_g)_{D(g)\leq u}\ \middle|\
 s_g\in A_g,\quad \rho_{h,g}(s_g)=s_h\text{ if }D(h)\leq D(g)\leq u
 \right\}.
\]
The index is the element $g$ of the ring, not the lattice element $D(g)$.
Thus no choice of a representative for a basic open is needed.
When $D(g)=D(h)$, compatibility relates the components through the
canonical mutually inverse restriction maps.

Ring operations are componentwise. Two sections are equal if all their
components are equal. For $v\leq u$, restriction forgets components;
identity and composition follow from componentwise equality.

\begin{proposition}\label{prop:basic}
Evaluation at $g$ is a natural ring isomorphism
\[
 \epsilon_g:\OO_A(D(g))\xrightarrow{\ \cong\ }A_g.
\]
Its inverse sends $x$ to $(\rho_{h,g}(x))_{D(h)\leq D(g)}$.
\end{proposition}
\begin{proof}
Composition of the $\rho$ maps makes this family compatible. One inverse
law is $\rho_{g,g}=1$; the other is compatibility of sections. For
$D(h)\leq D(g)$, naturality is the same compatibility equation:
\[
 \epsilon_h(\restr{s}{D(h)})=\rho_{h,g}(\epsilon_g(s)).
\]
\end{proof}

\begin{proposition}[Limit property]
A ring map $T\to\OO_A(u)$ is uniquely determined by compatible maps
$T\to A_g$ for all $D(g)\leq u$. Every such compatible family induces a
map, by sending $t$ to the family of its images.
\end{proposition}

\section{Finite gluing}

\begin{lemma}[Separation on a basic cover]\label{lem:sep}
Suppose $D(g)=\bigvee_{i=1}^nD(h_i)$. The maps
$\rho_{h_i,g}:A_g\to A_{h_i}$ are jointly injective.
\end{lemma}
\begin{proof}
It suffices to consider $x$ restricting to zero. Write $f_g(s)x=f_g(y)$
with $s$ a power of $g$. For each $i$, kernel detection gives a power
$t_i$ of $h_i$ with $t_iy=0$. Choose these finitely many witnesses.
Since $D(h_i)\leq D(t_i)$, the Nullstellensatz supplies
$g^p=\sum_i t_ic_i$. Hence $yg^p=0$, and cancellation of the units
$f_g(g^p)$ and $f_g(s)$ gives $x=0$.
\end{proof}

\begin{theorem}[Basic sheaf lemma]\label{thm:basic-sheaf}
Suppose $D(g)=\bigvee_iD(h_i)$ and $x_i\in A_{h_i}$ have equal images
in $A_{h_ih_j}$ for every $i,j$. There is a unique $x\in A_g$ whose
image in $A_{h_i}$ is $x_i$ for every $i$.
\end{theorem}
\begin{proof}
Uniqueness is Lemma~\ref{lem:sep}, so the type of amalgams is a
proposition and merely existing fraction data may be used.
Write $f_{h_i}(s_i)x_i=f_{h_i}(y_i)$, where $s_i$ is a power of $h_i$.
Agreement on overlaps implies that some power $(h_ih_j)^{m_{ij}}$
annihilates $y_is_j-y_js_i$. Choose all these finitely many exponents
and a common bound $M$. Put
\[
 w_i=h_i^{M+1},\qquad k_i=w_is_i,\qquad z_i=w_iy_i.
\]
Then $D(k_i)=D(h_i)$ and the cross relations hold in $A$ itself:
\[
 z_ik_j=z_jk_i.
\]
The positive exponent in $w_i$ ensures $D(k_i)=D(h_i)$ even when the
chosen denominator $s_i$ is the zeroth power.

By the Nullstellensatz, merely $g^p=\sum_i k_ic_i$. Set
$Y=\sum_i z_ic_i$. The cross relations give $Yk_j=z_jg^p$, so define
\[
 x=f_g(Y)\,f_g(g^p)^{-1}.
\]
In $A_{k_j}$, cancellation of $f_{k_j}(g^p)$ gives
$f_{k_j}(k_j)\rho_{k_j,g}(x)=f_{k_j}(z_j)$.
The image of $x_j$ satisfies the same equation, and $f_{k_j}(k_j)$ is a
unit. Thus their images agree in $A_{k_j}$. Finally
$A_{h_j}\to A_{k_j}$ is injective by Lemma~\ref{lem:sep}, since their
basic opens are equal. Therefore $\rho_{h_j,g}(x)=x_j$.
The proof includes $n=0$: in that case $D(g)=\bot$ and $A_g$ is the
zero ring.
\end{proof}

\begin{definition}
A presheaf of rings on a distributive lattice $L$ consists of rings
$\OO(u)$ and restriction maps $\OO(u)\to\OO(v)$ for $v\leq u$,
satisfying identity and composition. It is a \emph{sheaf} if every
compatible family on a finite cover $u=\bigvee_i u_i$ has a unique
amalgam in $\OO(u)$. Compatibility means equality after restriction to
$u_i\wedge u_j$. The empty cover is included, so $\OO(\bot)$ is the
terminal ring.
\end{definition}

\begin{theorem}\label{thm:sheaf}
The presheaf $\OO_A$ is a sheaf on $\Zar(A)$.
\end{theorem}
\begin{proof}
First establish separation on arbitrary finite covers. If two sections
agree on the basic generators of a finite presentation of $v$, their
components at any $D(h)\leq v$ agree: the opens $D(ha_j)$ cover $D(h)$,
so apply Lemma~\ref{lem:sep}. This extends agreement through binary and
finite joins, including bottom, and proves separation.

Let $u=\bigvee_i u_i$ and let $s_i\in\OO_A(u_i)$ be compatible. For
$D(g)\leq u$, consider sections $y\in\OO_A(D(g))$ agreeing with every
$s_i$ on $D(g)\wedge u_i$. There is at most one such section, by
separation on the cover $D(g)=\bigvee_i(D(g)\wedge u_i)$.

For existence, choose finite presentations of the $u_i$ and concatenate
them, retaining the chart index of each generator $a_j$. The opens
$D(ga_j)$ cover $D(g)$, since
\[
 D(g)\wedge\bigvee_jD(a_j)=\bigvee_jD(ga_j).
\]
Take the appropriate components of the $s_i$ on these opens.
Compatibility makes them agree on overlaps.
Theorem~\ref{thm:basic-sheaf} and Proposition~\ref{prop:basic} produce
$y$; separation verifies its agreement with each $s_i$ throughout
$D(g)\wedge u_i$. Since $y$ is uniquely characterised, this construction
is independent of the chosen presentations.

Denote this section by $y_g$ and set $s_g=\epsilon_g(y_g)$. If
$D(h)\leq D(g)$, the restriction of $y_g$ is the uniquely characterised
$y_h$. Naturality of evaluation gives $\rho_{h,g}(s_g)=s_h$.
Thus the $s_g$ form a section of $\OO_A(u)$, and its restrictions are
the $s_i$. Separation proves uniqueness. This argument applies unchanged
to the empty cover.
\end{proof}

\section{Extending maps from basic opens}

\begin{theorem}\label{thm:extension}
Let $G$ be a sheaf of rings on a distributive lattice $M$, let
$\pi:\Zar(A)\to M$ be a lattice map, and suppose given ring maps
\[
 \varphi_g:A_g\longrightarrow G(\pi D(g))
\]
commuting with restriction. There is a unique family of ring maps
\[
 \Phi_u:\OO_A(u)\longrightarrow G(\pi u)
\]
commuting with restriction and satisfying
$\Phi_{D(g)}=\varphi_g\epsilon_g$.
\end{theorem}
\begin{proof}
For $s\in\OO_A(u)$, characterise $t\in G(\pi u)$ by
\begin{equation}\label{eq:fits}
 \restr{t}{\pi D(g)}=\varphi_g(s_g)\qquad(D(g)\leq u).
\end{equation}
A finite presentation $u=\bigvee_jD(a_j)$ gives a cover
$\pi u=\bigvee_j\pi D(a_j)$. Separation in $G$ proves that at most one
$t$ satisfies \eqref{eq:fits}.

The elements $\varphi_{a_j}(s_{a_j})$ are compatible: their restrictions
to $\pi D(a_ja_k)=\pi D(a_j)\wedge\pi D(a_k)$ coincide by the assumed
compatibility of $\varphi$ and of $s$. Glue them to $t$.
For an arbitrary $D(g)\leq u$, the opens $\pi D(ga_j)$ cover
$\pi D(g)$. On each, both sides of \eqref{eq:fits} agree by restriction
compatibility. Separation in $G$ proves the equation. Thus $t$ is
independent of the presentation; put $\Phi_u(s)=t$.

The unit, sum and product of fitting elements fit the corresponding
unit, sum and product of sections. Uniqueness therefore makes $\Phi_u$
a ring map. The same argument proves compatibility with restriction and
the formula on basic opens. Any family with these properties satisfies
\eqref{eq:fits}, proving uniqueness.
\end{proof}

\section{Invertibility opens and local structure}

\begin{definition}
For a presheaf of rings $\OO$ on $L$, an \emph{invertibility open} of
$s\in\OO(u)$ is an element $v\leq u$ such that, for all $w\leq u$,
\[
 w\leq v\quad\Longleftrightarrow\quad \restr{s}{w}\in\units{\OO(w)}.
\]
An \emph{invertibility map} assigns such an open $D_u(s)$ to every
section.
\end{definition}

\begin{lemma}\label{lem:inv}
Invertibility opens are unique. If $v=D_u(s)$ and $w\leq u$, then
\[
 D_w(\restr{s}{w})=w\wedge v.
\]
For every invertibility map,
\[
 D_u(1)=u,\qquad D_u(st)=D_u(s)\wedge D_u(t).
\]
\end{lemma}
\begin{proof}
An invertibility open is the greatest open below $u$ on which $s$ is a
unit, so it is unique. Restriction follows by testing on opens below
$w$. The unit is invertible everywhere. In a commutative ring a product
is a unit exactly when both factors are units, which proves the product
formula by the same characterisation.
\end{proof}

\begin{lemma}\label{lem:local-units}
In a sheaf of rings, being a unit is local for finite covers. Consequently
the defining property of an invertibility open may be tested on a basis
of finite covers. If $u=\bigvee_i u_i$ and the restrictions of $s$ have
invertibility opens $v_i\leq u_i$, then $\bigvee_i v_i$ is the
invertibility open of $s$.
\end{lemma}
\begin{proof}
Local inverses agree on intersections by uniqueness of inverses, and
therefore glue. The product of the glued inverse with the original
section is one by separation. This also includes the empty cover.
For the basis criterion, present any test open as a finite join of basic
opens and apply this locality of units.

For the last assertion put $v=\bigvee_i v_i$. The section is a unit on
each $v_i$, hence on $v$. Conversely, if it is a unit on $w\leq u$,
then $w\wedge u_i\leq v_i$ for every $i$, and
$w=\bigvee_i(w\wedge u_i)\leq v$.
\end{proof}

\begin{definition}
A \emph{locally ringed lattice} is a distributive lattice $L$ with a
sheaf of rings $\OO$ and an invertibility map satisfying
\[
 D_u(0)=\bot,\qquad D_u(s+t)\leq D_u(s)\vee D_u(t).
\]
Together with Lemma~\ref{lem:inv}, these say that
$D_u:\OO(u)\to\mathord\downarrow u$ is a support.
\end{definition}

\begin{theorem}\label{thm:affine}
The ringed lattice $(\Zar(A),\OO_A)$ is locally ringed. For a section
$s\in\OO_A(D(g))$ with $\epsilon_g(s)=a/g^n$,
\begin{equation}\label{eq:basic-inv}
 D_{D(g)}(s)=D(ga).
\end{equation}
\end{theorem}
\begin{proof}
For $D(h)\leq D(g)$, the fraction equation shows that the image of $s$
in $A_h$ is a unit exactly when $f_h(a)$ is a unit. By
Proposition~\ref{prop:bridge}, this is equivalent to
$D(h)\leq D(a)$, hence to $D(h)\leq D(ga)$.
Evaluation is an isomorphism, so it preserves and reflects units.
The basis criterion proves \eqref{eq:basic-inv}. Uniqueness makes the
answer independent of the fraction. Finite presentations and
Lemma~\ref{lem:local-units} now give an invertibility open for every
section on every $u$.

For zero, a basic open $D(b)$ where zero is a unit satisfies
$D(b)\leq D(0)=\bot$ by the bridge. Finite generation therefore gives
$D_u(0)=\bot$.

For addition, test below $D_u(s+t)$ on a basic open $D(b)$ and write the
evaluations of the restrictions of $s,t$ as $a/p,c/q$ in $A_b$, with
$p,q$ powers of $b$. Since their sum is a unit, the identity
\[
 f_b(qa+pc)=f_b(p)f_b(q)\bigl(\epsilon_b(\restr{s}{D(b)})+
                              \epsilon_b(\restr{t}{D(b)})\bigr)
\]
shows that $f_b(qa+pc)$ is a unit. The bridge and the support laws of
$D$ give
\[
 D(b)\leq D(qa+pc)\leq D(a)\vee D(c).
\]
Intersecting with $D(b)$ gives
$D(b)\leq D(ba)\vee D(bc)$. By \eqref{eq:basic-inv} and restriction
compatibility these are contained respectively in $D_u(s)$ and
$D_u(t)$. Finite generation proves the required addition inequality.
\end{proof}

We denote this locally ringed lattice by $\Spec A$.

A geometric morphism $X\to Y$ of ringed lattices consists of a lattice
map $\pi:L_Y\to L_X$ and compatible ring maps
$\pi^\sharp_u:\OO_Y(u)\to\OO_X(\pi u)$. It is \emph{local} when
\[
 \pi(D^Y_u(s))=D^X_{\pi u}(\pi^\sharp_u(s)).
\]
The forward inequality always holds: ring maps preserve units and
$\pi^\sharp$ commutes with restriction. Thus locality is precisely the
reverse inequality. Notice the contravariance of the lattice map.

\section{The universal property of the affine spectrum}

Write $\Gamma X=\OO_X(\top)$ for the ring of global sections of a locally
ringed lattice $X$. For $a\in A$ and $u\in\Zar(A)$, the family
$(f_g(a))_{D(g)\leq u}$ is a section of $\OO_A(u)$; denote it by $a|_u$.
These sections commute with restriction and define ring maps
$A\to\OO_A(u)$.

\begin{proposition}\label{prop:global}
The homomorphism
\[
 \alpha_A:A\longrightarrow\Gamma(\Spec A),\qquad a\longmapsto a|_\top,
\]
is an isomorphism. Moreover,
\[
 D_\top(a|_\top)=D(a).
\]
\end{proposition}
\begin{proof}
Since $D(1)=\top$, evaluation identifies $\Gamma(\Spec A)$ with $A_1$.
The identity of $A$ inverts every power of $1$, so the universal property
of localisation supplies $b:A_1\to A$ with $bf_1=1_A$.
The maps $f_1b$ and $1_{A_1}$ agree on $A$ and hence are equal by
Corollary~\ref{cor:agree}. Thus $f_1$ is an isomorphism, and so is
$\alpha_A$. Explicitly, its inverse sends $s$ to $b(\epsilon_1(s))$;
the inverse laws follow from compatibility of the components of $s$.
The support formula is \eqref{eq:basic-inv} applied to $a/1$ on $D(1)$.
\end{proof}

\begin{theorem}[Universal property of $\Spec$]\label{thm:spec-universal}
For any locally ringed lattice $X$ and commutative ring $A$, local
morphisms $X\to\Spec A$ correspond bijectively to ring homomorphisms
$A\to\Gamma X$:
\[
 \Hom_{\mathrm{loc}}(X,\Spec A)
 \cong \Hom_{\mathrm{Ring}}(A,\Gamma X).
\]
The correspondence sends a morphism to its map on global sections,
composed with $\alpha_A$.
\end{theorem}
\begin{proof}
Let $X$ have lattice $M$, and let $\varphi:A\to\OO_X(\top)$ be a ring
homomorphism. We construct a local morphism with inverse-image map
$\pi:\Zar(A)\to M$ and maps
$\pi^\sharp_u:\OO_A(u)\to\OO_X(\pi u)$.

\emph{The lattice map.} The function
\[
 a\longmapsto D^X_\top(\varphi(a))
\]
is a support, because $\varphi$ is a ring map and $D^X_\top$ is a
support. Theorem~\ref{thm:support} supplies a unique lattice map $\pi$
with
\begin{equation}\label{eq:pi-D}
 \pi D(g)=D^X_\top(\varphi(g)).
\end{equation}

\emph{The maps on basic opens.} Restrict $\varphi$ to $\pi D(g)$.
The defining property of the invertibility open says that the image of
$g$ is a unit there, and therefore so is the image of every power of $g$.
The universal property of $A_g$ gives a unique homomorphism
\[
 \varphi_g:A_g\longrightarrow\OO_X(\pi D(g)),\qquad
 \varphi_g(f_g(a))=\restr{\varphi(a)}{\pi D(g)}.
\]
For $D(h)\leq D(g)$ the two maps
\[
 A_g\longrightarrow\OO_X(\pi D(h))
\]
obtained by restriction of $\varphi_g$ and by $\varphi_h\rho_{h,g}$
agree on $A$. They are equal by Corollary~\ref{cor:agree}.

\emph{The sheaf map.} Theorem~\ref{thm:extension} extends these maps
uniquely to a family $\pi^\sharp_u$ commuting with restriction, with
\begin{equation}\label{eq:sharp-basic}
 \pi^\sharp_{D(g)}(s)=\varphi_g(\epsilon_g(s)).
\end{equation}
This defines a morphism of ringed lattices. On global sections it sends
$a|_\top$ to $\varphi(a)$, by the defining equation of the extension at
$D(1)=\top$.

\emph{Locality on basic opens.} Suppose
$\epsilon_g(s)=a/g^n$, put $w=\pi D(g)$, and write
$t=\pi^\sharp_{D(g)}(s)$. The fraction equation gives
\[
 y t=\restr{\varphi(a)}{w},\qquad
 y=\restr{\varphi(g^n)}{w}\in\units{\OO_X(w)}.
\]
The product law for invertibility opens and restriction compatibility
give
\begin{align*}
 D^X_w(t)
 &=D^X_w(yt)
  =w\wedge D^X_\top(\varphi(a))\\
 &=\pi D(g)\wedge\pi D(a)
  =\pi D(ga)
  =\pi\bigl(D^{\Spec A}_{D(g)}(s)\bigr).
\end{align*}
Here the last equality is \eqref{eq:basic-inv}.

\emph{Locality on arbitrary opens.} Present $u$ as
$\bigvee_iD(a_i)$. Both $u$ and $\pi u$ have the corresponding finite
covers. Apply the join formula of Lemma~\ref{lem:local-units} to $s$
and $\pi^\sharp_u(s)$, using compatibility of $\pi^\sharp$ with
restriction and the basic-open calculation. It follows that
\[
 D^X_{\pi u}(\pi^\sharp_u(s))
 =\bigvee_i\pi\bigl(D^{\Spec A}_{D(a_i)}(\restr{s}{D(a_i)})\bigr)
 =\pi\bigl(D^{\Spec A}_u(s)\bigr).
\]
Thus the constructed morphism is local.

\emph{Uniqueness.} Let a local morphism with lattice map $\psi$ have
global component $\varphi$. Proposition~\ref{prop:global} and locality
force
\[
 \psi D(a)
 =\psi\bigl(D^{\Spec A}_\top(a|_\top)\bigr)
 =D^X_\top(\varphi(a)).
\]
Universal support gives $\psi=\pi$. With the lattice maps identified,
consider its map on a basic open, composed with the inverse of
$\epsilon_g$. This is a map $A_g\to\OO_X(\pi D(g))$ agreeing with
$\varphi_g$ on $A$, by restriction compatibility and the hypothesis on
global sections. Corollary~\ref{cor:agree} identifies it with
$\varphi_g$. The uniqueness clause of
Theorem~\ref{thm:extension} then identifies all the sheaf maps with
$\pi^\sharp$.

We have constructed a local morphism with prescribed global component
and proved it unique. Consequently, taking the global component and
constructing the morphism are mutually inverse operations.
\end{proof}

In particular, the identity homomorphism of $\Gamma X$ determines a
canonical local morphism
\[
 X\longrightarrow\Spec(\Gamma X).
\]
Its construction uses the invertibility support of $X$; $X$ need not
already be given with affine charts.

\section{Restriction, affineness and qcqs schemes}

\begin{proposition}\label{prop:restriction}
Let $X=(L,\OO,D)$ be a locally ringed lattice and $e\in L$. There is a
locally ringed lattice $X|_e$ on
\[
 \mathord\downarrow e=\{u\in L\mid u\leq e\},
\]
with section ring $\OO(u)$ at $u\leq e$ and the inherited restriction
maps. Its invertibility support is the support of $X$, viewed in this
principal ideal. The inclusion $X|_e\to X$ is local.
\end{proposition}
\begin{proof}
Meets, joins and bottom are inherited from $L$, and top is $e$.
A finite cover in $\mathord\downarrow e$ is a finite cover of the same
element in $L$, so the sheaf condition is inherited. For $s\in\OO(u)$,
the open $D_u(s)$ already lies below $u\leq e$. Testing it on opens
below $u$ gives exactly the same invertibility characterisation as in
$X$. Its support laws are inherited as well.

The inverse-image map of the inclusion is
$u\mapsto u\wedge e$, a lattice map from $L$ to
$\mathord\downarrow e$. Its maps on rings are restriction
$\OO(u)\to\OO(u\wedge e)$. Their naturality is composition of
restrictions. Locality is the identity
\[
 D_{u\wedge e}(\restr{s}{u\wedge e})=D_u(s)\wedge e,
\]
which follows from Lemma~\ref{lem:inv} and $D_u(s)\leq u$.
\end{proof}

\begin{proposition}[Universal property of restriction]\label{prop:slice-up}
A local morphism $f:Z\to X$ factors through $X|_e\to X$ if and only if
$f^*(e)=\top$. The factorisation is unique.
\end{proposition}
\begin{proof}
If $f^*(e)=\top$, restrict its lattice map to
$\mathord\downarrow e$ and retain its section maps there. This is a
lattice map because the new top $e$ maps to top; the other lattice laws
and locality are inherited. Composing with the inclusion recovers $f$:
on lattices, $f^*(u\wedge e)=f^*(u)$, and on rings the equality is
naturality with respect to restriction. Conversely, the inverse image
of $e$ under the inclusion is top in $\mathord\downarrow e$.
Uniqueness follows because every $u\leq e$ satisfies $u\wedge e=u$,
and restriction between these equal opens is the identity.
\end{proof}

An isomorphism of ringed lattices may be characterised by a bijective
inverse-image map on lattices and bijective maps on all the section
rings. The inverse lattice map and inverse ring maps supply its inverse
morphism; naturality follows from that of the original maps.

\begin{lemma}\label{lem:iso-local}
Every isomorphism of ringed lattices between locally ringed lattices is
local.
\end{lemma}
\begin{proof}
Write its lattice map as $\pi:L_Y\to L_X$. Given $s\in\OO_Y(u)$, put
$v=D^X_{\pi u}(\pi^\sharp_u(s))$ and let $w=\pi^{-1}(v)$. Since
$v\leq\pi u$ and $\pi$ is a lattice isomorphism, $w\leq u$.
Naturality shows that the image of $\restr{s}{w}$ under the section-ring
isomorphism at $w$ is a unit. Ring isomorphisms reflect units, so
$\restr{s}{w}$ is a unit. Hence $w\leq D^Y_u(s)$ and
$v\leq\pi D^Y_u(s)$. The opposite inequality holds for every morphism
of ringed lattices, proving locality.
\end{proof}

Identities and composites of local morphisms are local, directly from
the defining support equation. Isomorphisms compose. Restriction between
equal opens is a ring isomorphism, since the two restriction maps are
inverse by the restriction laws. Consequently, identifying two lattice
maps pointwise and identifying their section maps through these
restriction isomorphisms preserves the property of being an isomorphism.

\begin{proposition}\label{prop:spec-map}
A ring homomorphism $\beta:B\to A$ induces a local morphism
\[
 \Spec\beta:\Spec A\longrightarrow\Spec B.
\]
It is an isomorphism if $\beta$ is a ring isomorphism.
\end{proposition}
\begin{proof}
Apply Theorem~\ref{thm:spec-universal} to
$\alpha_A\beta:B\to\Gamma(\Spec A)$. If
$\psi:X\to\Spec A$ has global component $\varphi:A\to\Gamma X$,
the composite $(\Spec\beta)\psi$ has global component $\varphi\beta$,
by composition of the maps on global sections.

If $\gamma:A\to B$ is inverse to $\beta$, the two composites of
$\Spec\beta$ and $\Spec\gamma$ have the same global components as
the respective identity morphisms. Uniqueness in
Theorem~\ref{thm:spec-universal} identifies them with those identities.
Thus the two induced morphisms are inverse isomorphisms.
\end{proof}

\begin{definition}
A locally ringed lattice $X$ is \emph{affine} if its canonical morphism
\[
 c_X:X\longrightarrow\Spec(\Gamma X),
\]
corresponding to the identity of $\Gamma X$, is an isomorphism.
\end{definition}

\begin{theorem}[Affine criterion]\label{thm:affine-criterion}
A locally ringed lattice is affine if and only if it is isomorphic to
$\Spec A$ for some commutative ring $A$.
\end{theorem}
\begin{proof}
If $X$ is affine, take $A=\Gamma X$. Conversely, let
$\psi:X\xrightarrow{\sim}\Spec A$ be an isomorphism. It is local by
Lemma~\ref{lem:iso-local}. Its global component is a ring isomorphism
\[
 \alpha:A\xrightarrow{\alpha_A}\Gamma(\Spec A)
             \xrightarrow{\Gamma\psi}\Gamma X.
\]
The second arrow is an isomorphism because the section map at top is
bijective and restriction between the equal top opens is bijective.
Proposition~\ref{prop:spec-map} gives the isomorphism
\[
 \Spec(\alpha^{-1}):\Spec A\xrightarrow{\sim}\Spec(\Gamma X).
\]
The composite $\Spec(\alpha^{-1})\psi$ is local and has global component
$\alpha\alpha^{-1}=1_{\Gamma X}$. By
Theorem~\ref{thm:spec-universal} it equals $c_X$. Therefore $c_X$ is an
isomorphism.
\end{proof}

\begin{definition}\label{def:qcqs}
An open $u$ of a locally ringed lattice $X$ is \emph{affine} if $X|_u$
is affine. An \emph{affine cover} is a finite family of affine opens
$u_1,\ldots,u_n$ with $\bigvee_i u_i=\top$. A \emph{qcqs scheme}, in
the lattice formulation, is a locally ringed lattice for which an affine
cover merely exists.
\end{definition}

The affine cover is a witness of a property, not part of the structure.
Its intersections are not required to be affine. In particular, two
affine opens with join top suffice to exhibit a qcqs scheme. The finite
cover may be empty.

\begin{theorem}[Principal affine opens]\label{thm:spec-basic}
For $g\in A$, there is an isomorphism of locally ringed lattices
\[
 (\Spec A)|_{D(g)}\cong\Spec(A_g).
\]
In particular, every principal open of an affine spectrum is affine.
\end{theorem}
\begin{proof}
Write $\beta:A\to A_g$ for localisation. The morphism
$\Spec\beta:\Spec(A_g)\to\Spec A$ sends $D(g)$ to top, since
$\beta(g)$ is a unit. Proposition~\ref{prop:slice-up} therefore gives
\[
 j:\Spec(A_g)\longrightarrow(\Spec A)|_{D(g)}.
\]
The inverse of evaluation in Proposition~\ref{prop:basic} is a ring map
\[
 A_g\longrightarrow\OO_A(D(g))
                =\Gamma((\Spec A)|_{D(g)}).
\]
Theorem~\ref{thm:spec-universal} gives a morphism $k$ in the other
direction. To show $kj=1$, compare their global components. It suffices
to check them on $\beta(a)$ by uniqueness for localisation; both give
the global section of $\beta(a)$. The universal property of $\Spec$
then proves equality of the morphisms.

For $jk=1$, compose both sides with the inclusion into $\Spec A$.
Their global components agree on every $a\in A$, so
Theorem~\ref{thm:spec-universal} gives equality after inclusion.
Uniqueness in Proposition~\ref{prop:slice-up} cancels the inclusion.
Thus $j$ and $k$ are inverse. Affineness follows from
Theorem~\ref{thm:affine-criterion}.
\end{proof}

\section{Doubling along an open}

Let $X=(L,\OO,D)$ be a locally ringed lattice and let $d\in L$.

\begin{definition}
The \emph{doubled lattice along $d$} is
\[
 L_2=\{(u,v)\in L\times L\mid u\wedge d=v\wedge d\}.
\]
For $x=(u,v)$ write $m(x)=u\wedge d=v\wedge d$.
\end{definition}

Both bounds and both operations are componentwise. Distributivity shows
that the operations preserve the defining equation. Order is
componentwise, and the two projections preserve finite joins and meets.

Recall the ring fibre product
\[
 A\times_C B=\{(a,b)\mid f(a)=g(b)\}
\]
for maps $f:A\to C$, $g:B\to C$. It is characterised by the fact that
two maps into $A,B$ with equal composites into $C$ induce a unique map
into $A\times_C B$. A pair is a unit if and only if its components are:
the images of their inverses coincide by uniqueness of inverses in $C$.

Define a presheaf on $L_2$ by
\[
 \OO_2(u,v)=\OO(u)\times_{\OO(m(u,v))}\OO(v).
\]
Restrictions are componentwise. The universal property of the fibre
product constructs these ring maps; their identity and composition laws
follow by the two projections.

\begin{theorem}\label{thm:double}
The presheaf $\OO_2$ is a sheaf, and
\[
 D^2_{(u,v)}(s,t)=\bigl(D_u(s),D_v(t)\bigr)
\]
is an invertibility support. Thus $(L_2,\OO_2,D^2)$ is a locally ringed
lattice.
\end{theorem}
\begin{proof}
For a compatible family on a finite cover of $(u,v)$, glue its first
components in $\OO(u)$ and its second components in $\OO(v)$, giving
$s,t$. Their restrictions to $m(u,v)$ agree: the corresponding overlap
opens of the covering family cover $m(u,v)$, by distributivity, and
the restrictions agree on each by the fibre-product conditions.
Separation gives agreement on all of $m(u,v)$. Thus $(s,t)$ is an
amalgam; its uniqueness follows componentwise. The argument includes
the empty cover.

The proposed invertibility open belongs to $L_2$. Indeed, since
$D_u(s)\leq u$,
\[
 D_u(s)\wedge d=m(u,v)\wedge D_u(s)
                 =D_{m(u,v)}(\restr{s}{m(u,v)}).
\]
The same expression for $t$ is equal, since the overlap sections agree.
For $(w,z)\leq(u,v)$, the restriction of $(s,t)$ is a unit if and only
if both component restrictions are units, equivalently
$w\leq D_u(s)$ and $z\leq D_v(t)$. This proves the invertibility
characterisation. The zero and addition support laws hold componentwise.
\end{proof}

\begin{proposition}[The two charts]\label{prop:charts}
Put $e_0=(\top,d)$ and $e_1=(d,\top)$. Then
\[
 e_0\vee e_1=\top,\qquad e_0\wedge e_1=(d,d).
\]
The original locally ringed lattice is isomorphic to each restriction
$X_2|_{e_i}$, by
\[
 \iota_0(u)=(u,u\wedge d),\qquad
 \iota_1(u)=(u\wedge d,u).
\]
More precisely, $\iota_i$ identifies $L$ with the opens below $e_i$,
and there are natural ring isomorphisms
\[
 \OO(u)\cong\OO_2(\iota_i(u))
\]
carrying invertibility opens to invertibility opens.
\end{proposition}
\begin{proof}
If $(a,b)\leq e_0$, then $b\leq d$ and
$b=b\wedge d=a\wedge d$. Thus $(a,b)=\iota_0(a)$. This proves the first
lattice identification, and the second is symmetric. Under the first,
the section isomorphism is
\[
 s\longmapsto(s,\restr{s}{u\wedge d}),
\]
with inverse first projection. Its inverse laws follow from the overlap
condition. It commutes with restriction, and
\[
 D^2_{\iota_0(u)}(s,\restr{s}{u\wedge d})
  =(D_u(s),D_u(s)\wedge d)=\iota_0(D_u(s)).
\]
For the second chart use $s\mapsto(\restr{s}{u\wedge d},s)$ and second
projection. Its overlap is $(u\wedge d)\wedge d$, equal to $u\wedge d$;
the restriction identities supply the same inverse laws.
\end{proof}

\subsection{Maps into a doubling}

There is a local \emph{fold morphism} $p:X_2\to X$. Its lattice map is
$a\mapsto(a,a)$, and its section maps are $s\mapsto(s,s)$.
Its locality follows from the componentwise invertibility support.

\begin{theorem}\label{thm:double-up}
For any locally ringed lattice $Z$, local morphisms $Z\to X_2$
correspond to triples $(f,u_0,u_1)$ with
\[
 f:Z\to X\text{ local},\qquad
 u_0\vee u_1=\top,\qquad u_0\wedge u_1=f^*(d).
\]
The correspondence takes $\psi$ to
$(p\psi,\psi^*e_0,\psi^*e_1)$. Thus the extents are the full inverse
images of the two charts.
\end{theorem}
\begin{proof}
For a triple as stated, define the inverse-image map by
\[
 \pi(a,b)=(u_0\wedge f^*a)\vee(u_1\wedge f^*b).
\]
It preserves bounds and joins. For meets, expand into four terms.
The two mixed terms lie below $u_0\wedge u_1=f^*d$, and
$a\wedge d=b\wedge d$ identifies the two coordinates there. Each
mixed term therefore lies below the required unmixed terms. This proves
preservation of meets.

For $(s,t)\in\OO_2(a,b)$, apply the section maps of $f$ and restrict to
$u_0\wedge f^*a$ and $u_1\wedge f^*b$. On their intersection the
images agree, by the defining overlap equation of $(s,t)$. Glue in
$Z$. Uniqueness of the amalgam proves preservation of ring operations
and compatibility with restriction. The join formula for invertibility
opens, and locality of $f$, prove locality of the constructed map.

Its composite with the fold has lattice map $f^*$ and the same section
maps as $f$, by separation. Its chart inverse images are $u_0,u_1$,
using $f^*d\leq u_0,u_1$.

Conversely every $(a,b)\in L_2$ decomposes as
\[
 (a,b)=\bigl(e_0\wedge(a,a)\bigr)
          \vee\bigl(e_1\wedge(b,b)\bigr).
\]
Indeed, its first coordinate is $a\vee(d\wedge b)=a$ and its second
is $(d\wedge a)\vee b=b$. Applying a lattice map to this identity
recovers the displayed formula for $\pi$. The section map is recovered
by separation on the same two pieces. These give both inverse laws.
\end{proof}

\section{The line with doubled origin}

Let $k$ be any commutative ring, put $A=k[X]$, and take
$d=D_A(X)$. Apply Theorem~\ref{thm:double} to $\Spec A$. The resulting
locally ringed lattice has compact opens
\[
 \left\{(u,v)\in\Zar(k[X])^2\ \middle|\
 u\wedge D(X)=v\wedge D(X)\right\}
\]
and section rings
\[
 \OO_2(u,v)=
 \OO_A(u)\times_{\OO_A(u\wedge D(X))}\OO_A(v).
\]
By Proposition~\ref{prop:charts}, its two chart opens
$(\top,D(X))$ and $(D(X),\top)$ cover it, and each chart is a copy of
$\Spec k[X]$. Their intersection corresponds in either chart to
$D(X)$, whose section ring is $k[X,X^{-1}]$ by
Proposition~\ref{prop:basic}. The transition identifies the two
coordinates, rather than inverting them. This is the line with doubled
origin: two copies of the affine line identified along the punctured
line. The construction permits an arbitrary base ring, including the
zero ring.

\begin{theorem}
The line with doubled origin over any commutative ring $k$ is a qcqs
scheme.
\end{theorem}
\begin{proof}
Each chart restriction is isomorphic to $\Spec k[X]$ by
Proposition~\ref{prop:charts}, now regarded as an isomorphism with the
restriction of Proposition~\ref{prop:restriction}. Both restrictions
are affine by Theorem~\ref{thm:affine-criterion}. The chart opens
$e_0=(\top,D(X))$ and $e_1=(D(X),\top)$ satisfy
\[
 e_0\vee e_1=(\top\vee D(X),D(X)\vee\top)=(\top,\top).
\]
They therefore form a two-element affine cover, giving the property in
Definition~\ref{def:qcqs}. No nontriviality hypothesis on $k$ is needed.
\end{proof}

\subsection{Global sections}

Continue to write $A=k[X]$, and denote the doubled-origin line by $S$.

\begin{lemma}\label{lem:punctured-injective}
The localisation homomorphism $A\to A_X$ is injective. Consequently,
global sections of $\Spec A$ which agree on $D(X)$ are equal.
\end{lemma}
\begin{proof}
Multiplication by $X$ is injective on $k[X]$: equality of coefficients
in $Xp=Xq$ gives equality of the coefficients of $p$ and $q$. Induction
shows that multiplication by every power $X^n$ is injective. If $a,b$
have the same image in $A_X$, kernel detection for localisation gives
some $n$ with $X^n(a-b)=0$. Hence $a=b$.

By Proposition~\ref{prop:global}, write the two global sections as
$a|_\top,b|_\top$. Their components on $D(X)$ are $f_X(a),f_X(b)$.
Agreement on $D(X)$ therefore implies $a=b$, proving the second claim.
\end{proof}

\begin{proposition}\label{prop:double-global}
Every global section $(s,t)$ of $S$ has $s=t$. The diagonal homomorphism
\[
 \delta:k[X]\longrightarrow\Gamma S,\qquad
 a\longmapsto(a|_\top,a|_\top),
\]
is a ring isomorphism.
\end{proposition}
\begin{proof}
The two components of a global section agree on
$\top\wedge D(X)=D(X)$ by the fibre-product condition. They are equal by
Lemma~\ref{lem:punctured-injective}. The diagonal map is a ring map by
the fibre-product universal property. An inverse sends $(s,t)$ to the
element of $k[X]$ corresponding to its first component under
Proposition~\ref{prop:global}. One inverse law is the corresponding law
for $\Gamma(\Spec A)\cong A$; the other uses $s=t$.
\end{proof}

Equivalently, under the affine global-section identifications,
\[
 \Gamma S\cong k[X]\times_{k[X,X^{-1}]}k[X]\cong k[X].
\]
The two chart restrictions of a global function thus give the same
polynomial. No assumption that $k$ is a domain is involved.

\subsection{Non-affineness}

\begin{theorem}\label{thm:double-not-affine}
If the doubled-origin line over $k$ is affine, then $1=0$ in $k$.
In particular, if $1\ne0$ in $k$, it is a non-affine qcqs scheme.
\end{theorem}
\begin{proof}
Let $\pi:\Zar(\Gamma S)\to L_2$ be the inverse-image lattice map of
the canonical morphism $c_S:S\to\Spec(\Gamma S)$. Call $(u,v)\in L_2$
\emph{diagonal} when $u=v$. Diagonal opens are closed under finite joins,
including the empty join, since the lattice operations are componentwise.

For a global section $r=(s,t)$, Proposition~\ref{prop:double-global}
gives $s=t$. Its invertibility open is consequently diagonal:
\[
 D^S_\top(r)=(D^A_\top(s),D^A_\top(t)).
\]
By the construction of the canonical morphism,
\[
 \pi D_{\Gamma S}(r)=D^S_\top(r).
\]
Every element of $\Zar(\Gamma S)$ merely has a finite presentation by
generators $D_{\Gamma S}(r_i)$. Since $\pi$ preserves finite joins,
every open in its image is diagonal. This argument does not require an
identification of $\Zar(\Gamma S)$ with $\Zar(k[X])$.

If $S$ is affine, $c_S$ is an isomorphism, so $\pi$ is surjective. In
particular the first chart $e_0=(\top,D_A(X))$ lies in its image. It must
therefore be diagonal, giving $D_A(1)=\top=D_A(X)$. The localisation
bridge, or the formal Nullstellensatz, now supplies $n$ and $q\in k[X]$
with
\[
 1=1^n=Xq.
\]
Apply the ring homomorphism $k[X]\to k$ evaluating $X$ at zero. It gives
$1=0\,q(0)=0$ in $k$, as required.
\end{proof}

\subsection{Points and their extents}

The inclusion $k\to k[X]$ followed by $\delta$ of
Proposition~\ref{prop:double-global} defines a morphism $S\to\Spec k$
by Theorem~\ref{thm:spec-universal}. For a $k$-algebra $B$, a
$B$-valued point over $k$ means a local morphism $\Spec B\to S$ whose
composite to $\Spec k$ is induced by the given structure map $k\to B$.

\begin{theorem}\label{thm:double-points}
For every $k$-algebra $B$ there is an equivalence
\[
 \Hom_k(\Spec B,S)\cong
 \left\{(t,u_0,u_1)\ \middle|\
 \begin{array}{l}
 t\in B,\quad u_0,u_1\in\Zar(B),\\
 u_0\vee u_1=\top,\quad u_0\wedge u_1=D_B(t)
 \end{array}\right\}.
\]
Here $u_i$ is the full inverse image of the $i$th chart. For a
$k$-algebra map $h:B\to B'$, precomposition with $\Spec h$ sends the
coordinates to
\[
 (h(t),\Zar(h)(u_0),\Zar(h)(u_1)),
\]
where $\Zar(h)(D_B(b))=D_{B'}(h(b))$.
\end{theorem}
\begin{proof}
Apply Theorem~\ref{thm:double-up} with $Z=\Spec B$ and
$X=\Spec k[X]$. A map $f:\Spec B\to\Spec k[X]$ over $k$ is, by
Theorem~\ref{thm:spec-universal} and Proposition~\ref{prop:global}, a
$k$-algebra homomorphism $k[X]\to B$. The polynomial universal property
identifies it with the image $t$ of $X$. Locality gives
$f^*D(X)=D_B(t)$. This identifies the triples and proves the equivalence.

For base change, the global component of the composite is the composite
algebra homomorphism, so the coordinate becomes $h(t)$. The inverse
images of the chart opens compose, giving the two asserted extents.
The generator formula for $\Zar(h)$ follows from the support equation
for the spectrum morphism and Proposition~\ref{prop:global}.
\end{proof}

\subsection{Complementary opens and idempotents}

Write $\Idem(C)=\{e\in C\mid e^2=e\}$. The next two results explain
how the pair of chart extents becomes an idempotent modulo $t$.

\begin{proposition}\label{prop:idempotents}
In any locally ringed lattice $X$, idempotents of $\Gamma X$ correspond
to complementary pairs of opens. The correspondence is
\[
 e\longmapsto\bigl(D_\top(e),D_\top(1-e)\bigr).
\]
\end{proposition}
\begin{proof}
For $e^2=e$, the identities $e(1-e)=0$ and $e+(1-e)=1$ show that the
two supports have meet bottom and join top.

Conversely let $u\vee v=\top$ and $u\wedge v=\bot$. Glue $1$ on $u$
with $0$ on $v$. They agree on the intersection because the ring on
bottom is terminal, by the empty-cover sheaf condition. Denote the
amalgam by $e$. Separation gives $e^2=e$.

More generally, a section with restrictions $1$ on $u$ and $0$ on $v$
has support $u$: restriction compatibility gives
$u\wedge D_\top(e)=u$ and $v\wedge D_\top(e)=\bot$, and distributivity
over $u\vee v=\top$ gives $D_\top(e)=u$. The same argument gives
$D_\top(1-e)=v$.

For the other inverse law, an idempotent unit is one, so
$\restr{e}{D_\top(e)}=1$. On $D_\top(1-e)$ the element $1-e$ is a
unit annihilating $e$, so the restriction of $e$ is zero. Separation
therefore recovers $e$ from its complementary pair.
\end{proof}

\begin{proposition}\label{prop:quotient-lattice}
For $t\in B$, there is a canonical lattice isomorphism
\[
 \Zar(B/(t))\cong[D_B(t),\top].
\]
If $q:B\to B/(t)$ is the quotient map, the image of $D(q(b))$ is
$D_B(t)\vee D_B(b)$. The interval has bottom $D_B(t)$.
\end{proposition}
\begin{proof}
Both lattices have the universal property of a support of $B$ killing
$t$. For the interval use the support
$b\mapsto D_B(t)\vee D_B(b)$. A support $d:B\to N$ with $d(t)=\bot$
extends uniquely to $\Zar(B)$; restriction of that lattice map to the
interval preserves bottom and has the desired factorisation. Uniqueness
follows by precomposing with $u\mapsto D_B(t)\vee u$, which is onto the
interval.

For the quotient, a support killing $t$ descends uniquely to $B/(t)$.
Indeed, if $b'=b+tc$, then $d(b')\leq d(b)\vee d(tc)=d(b)$, and the
reverse inequality follows in the same way. The support laws descend
on representatives. Universal support for $B/(t)$ now gives the second
universal property. The two resulting lattice maps are inverse by
uniqueness. Their values on $q(b)$ give the formula.
\end{proof}

\begin{theorem}\label{thm:double-idempotents}
For every $k$-algebra $B$, naturally in $B$,
\[
 \Hom_k(\Spec B,S)\cong\sum_{t\in B}\Idem(B/(t)).
\]
\end{theorem}
\begin{proof}
For fixed $t$, the pairs of Theorem~\ref{thm:double-points} are precisely
the complementary pairs in the interval $[D_B(t),\top]$.
Proposition~\ref{prop:quotient-lattice},
Proposition~\ref{prop:idempotents} applied to $\Spec(B/(t))$, and
Proposition~\ref{prop:global} identify them with $\Idem(B/(t))$.

Explicitly, if an idempotent is represented by $q(b)$, its extents are
\[
 u_0=D_B(t)\vee D_B(b),\qquad
 u_1=D_B(t)\vee D_B(1-b).
\]
This is independent of the representative by the quotient-lattice
isomorphism. A map $h:B\to B'$ induces
$\bar h:B/(t)\to B'/(h(t))$, and the formula commutes with the action of
$\Zar(h)$. Thus base change sends $(t,e)$ to $(h(t),\bar h(e))$.
The formula includes the zero ring: if $t$ is a unit, $B/(t)$ has just
one idempotent and the extents are $(\top,\top)$.
\end{proof}

\section{Gluing with a transition}

Let $X_1,X_2,Y$ be locally ringed lattices and let
$a_i:Y\to X_i$ be local morphisms, in geometric notation. Write
$p_i:L_i\to M$ for their inverse-image lattice maps. There is a
componentwise construction on the pullback lattice
\[
 L=\{(u,v)\in L_1\times L_2\mid p_1(u)=p_2(v)\}.
\]
At $(u,v)$, put $w=p_1(u)=p_2(v)$ and define
\[
 \OO(u,v)=\OO_1(u)\times_{\OO_Y(w)}\OO_2(v),
\]
using the two section maps of $a_i$. Operations and restrictions are
componentwise. As in Theorem~\ref{thm:double}, compatible sections glue
in $X_1$ and $X_2$, and separation in $Y$ verifies the overlap equation
for the amalgams. Locality of $a_i$ ensures that
\[
 D_{(u,v)}(s,t)=(D^1_u(s),D^2_v(t))
\]
belongs to the pullback lattice and is the invertibility open.
The support laws hold componentwise. This construction gives a locally
ringed lattice for any such span; the chart statements below use the
particular open inclusions and transition isomorphism.

Now take $X_1=X_2=X$, $Y=X|_d$, and let
$\theta:X|_d\xrightarrow{\sim}X|_d$ be an automorphism. Take $a_1$ to
be the inclusion and $a_2$ the inclusion after $\theta$. Write the
result as $X^\theta$. Its opens and sections are
\begin{align*}
 L_\theta&=\{(u,v)\mid u\wedge d=\theta^*(v\wedge d)\},\\
 \OO_\theta(u,v)&=\{(s,t)\mid
 \restr{s}{u\wedge d}
 =\theta^\sharp_{v\wedge d}(\restr{t}{v\wedge d})\}.
\end{align*}
Equal opens in the section equation are identified. For the identity
transition this is the doubling construction.

\begin{proposition}\label{prop:transition-charts}
The opens $e_0=(\top,d)$ and $e_1=(d,\top)$ cover $X^\theta$, and
both chart restrictions are isomorphic to $X$.
\end{proposition}
\begin{proof}
An automorphism of $\mathord\downarrow d$ preserves its top $d$, so
the chart pairs belong to $L_\theta$, and their join is top.
Below $e_0$, the second component lies below $d$ and is uniquely
determined by the first:
\[
 v=(\theta^*)^{-1}(u\wedge d).
\]
The first projection identifies the chart lattice with $L_X$. On
sections the second leg into the overlap ring is an isomorphism, so
first projection from the fibre product is an isomorphism. These
projections commute with restriction and hence give a ringed-lattice
isomorphism; it is local by Lemma~\ref{lem:iso-local}.

Below $e_1$, one has $u=\theta^*(v\wedge d)$. Use second projection
instead: the first leg into the overlap ring is restriction between
equal opens and hence an isomorphism. This proves the second chart
identification. The second argument in fact needs no invertibility of
$\theta$.
\end{proof}

The construction and chart identifications do not assert a general
mapping property for gluing with a transition. Unlike the identity
doubling, a nontrivial transition need not admit a coordinate fold back
to $X$.

\section{The projective line}

\subsection{Inversion on the punctured chart}
Put $A=k[X]$, $L=A_X$, $f:A\to L$, and $x=f(X)$.
Let $x^{-1}$ be its inverse. The polynomial universal property gives a
homomorphism $\tau:A\to L$ fixing the image of $k$ and sending $X$ to
$x^{-1}$. Since $\tau(X)$ is a unit, localisation extends it uniquely to
\[
 \sigma:L\longrightarrow L,\qquad \sigma f=\tau.
\]
It fixes $k$ and satisfies $\sigma(x)=x^{-1}$. Ring maps preserve
inverses, so $\sigma(x^{-1})=x$. Thus $\sigma^2f$ and $f$ agree on $k$
and on $X$, and are equal by the polynomial universal property.
Corollary~\ref{cor:agree} gives $\sigma^2=1_L$.

Apply $\Spec$ and use Theorem~\ref{thm:spec-basic} to obtain the
automorphism
\[
 \theta:(\Spec k[X])|_{D(X)}\xrightarrow{\sim}
             (\Spec k[X])|_{D(X)}
\]
whose action on the Laurent ring is inversion. Define
$\mathbb P^1_k=(\Spec k[X])^\theta$ by the preceding construction.
Its two affine coordinates obey $X_1=X_0^{-1}$ on the overlap.
Proposition~\ref{prop:transition-charts} and
Theorem~\ref{thm:affine-criterion} give two affine chart restrictions
covering it. Hence $\mathbb P^1_k$ is a qcqs scheme.

\subsection{Global functions}

\begin{lemma}\label{lem:inversion-overlap}
If $p,q\in k[X]$ satisfy $f(p)=\sigma(f(q))$ in $L$, then $p=q$ and
both are constant polynomials.
\end{lemma}
\begin{proof}
Use induction on a finite coefficient list representing $q$, with the
statement quantified over all $p$. The conclusion is a property, so a
mere coefficient representation suffices. If $q=0$, injectivity of $f$
gives $p=0$.

Write $q=c+Xq'$. The overlap equation becomes
\[
 f(p)=f(c)+x^{-1}\tau(q').
\]
For $p'=X(p-c)$ it follows that $f(p')=\tau(q')=\sigma(f(q'))$.
By induction $p'=q'$ is constant. The constant coefficient of $p'$ is
zero, so $p'=q'=0$. Regularity of $X$ gives $p=c$, and consequently
$q=c+Xq'=c$ as well. This argument needs neither a degree function nor
a domain hypothesis.
\end{proof}

\begin{theorem}\label{thm:p1-global}
For every commutative ring $k$, the constant-section homomorphism
\[
 k\longrightarrow\Gamma(\mathbb P^1_k),\qquad
 c\longmapsto(c|_\top,c|_\top)
\]
is an isomorphism. Its inverse takes the constant coefficient of the
polynomial on the first chart.
\end{theorem}
\begin{proof}
Identify global sections on each affine chart with $k[X]$. Under
evaluation on the punctured chart, the overlap equation for a pair
$(p,q)$ is exactly
\[
 f(p)=\sigma(f(q)).
\]
Indeed, the overlap automorphism is obtained from $\Spec\sigma$ through
the principal-open isomorphism, so its action on sections is $\sigma$.
Functoriality of global sections identifies the transition equation
with the displayed one. Restriction between the equal overlap opens
is invertible, so the implication works in both directions.

Every pair of identical constants satisfies the equation because
$\sigma$ fixes $k$. Lemma~\ref{lem:inversion-overlap} shows conversely
that every global section is such a pair. Reading the constant
coefficient gives both inverse laws. The map is a ring homomorphism
by the fibre-product universal property.
\end{proof}

\subsection{Non-affineness and the structure map}

\begin{theorem}\label{thm:p1-not-affine}
If $\mathbb P^1_k$ is affine, then $1=0$ in $k$. Hence it is non-affine
whenever $1\ne0$ in $k$.
\end{theorem}
\begin{proof}
The two components of every global section are the same constant
polynomial, by Theorem~\ref{thm:p1-global}. Its invertibility open is
therefore diagonal, since the glued support is componentwise.
Let
\[
 \pi:\Zar(\Gamma\mathbb P^1_k)\longrightarrow L_\theta
\]
be the lattice map of the canonical affine morphism. Its generators map
to invertibility opens of global sections, so their images are diagonal.
Finite generation and preservation of joins imply that every image is
diagonal. Here only diagonal pairs that belong to $L_\theta$ are used;
their finite joins remain diagonal.

If the canonical map were an isomorphism, the chart
$(\top,D_A(X))$ would lie in its image and hence be diagonal. Thus
$D_A(X)=\top$. As in Theorem~\ref{thm:double-not-affine}, the
Nullstellensatz gives $1=Xq$, and evaluation at zero gives $1=0$ in
$k$. No identification of $\Zar(\Gamma\mathbb P^1_k)$ with $\Zar(k)$
is needed.
\end{proof}

\begin{proposition}\label{prop:p1-structure}
There is a canonical local morphism
\[
 \mathbb P^1_k\longrightarrow\Spec k
\]
whose restriction to either affine chart is $\Spec(k\to k[X])$.
\end{proposition}
\begin{proof}
Apply Theorem~\ref{thm:spec-universal} to the constant-section map
$k\to\Gamma(\mathbb P^1_k)$ of Theorem~\ref{thm:p1-global}.
For either chart, compose with its inclusion and its isomorphism from
$\Spec k[X]$. Its global component sends $c\in k$ to the global
section of the constant polynomial $c$: restrict the constant pair to
the chart, then take the corresponding component. This is also the
global component of $\Spec(k\to k[X])$. Uniqueness in
Theorem~\ref{thm:spec-universal} proves the chart compatibilities.
\end{proof}

\section{What has been formalised}

The development in \texttt{unit/sag/qcqs} checks in Unit: the entry point
\texttt{All.seg} reports 100 modules and no failures at the stage described
here. Its parameters are univalence and replacement with the associated
surjectivity principle. Function extensionality, propositional truncation
and effective small set quotients are derived. The development uses none
of the geometric axioms Loc, synthetic quasi-coherence or Zariski-local
choice.

The main mathematical correspondence is as follows. File names refer to
the directory above.

\small
\begin{longtable}{@{}>{\raggedright\arraybackslash}p{0.43\textwidth}>{\raggedright\arraybackslash}p{0.53\textwidth}@{}}
\textbf{Mathematics}&\textbf{Files and principal declarations}\\\hline
Localisation and the elementary characterisation&
\texttt{Localisation}, \texttt{LocalisationFractions}:
\texttt{IsLoc}, \texttt{Elementary}, \texttt{Agree}\\[4pt]
Radicals and the Zariski lattice&
\texttt{Radical}, \texttt{FamilyLattice}, \texttt{ZariskiLattice}:
\texttt{order}, \texttt{reflect}, \texttt{generated}, \texttt{presented}\\[4pt]
Universal support&
\texttt{DistLattice}, \texttt{ZariskiUniversal}:
\texttt{factor}, \texttt{factor-D}, \texttt{unique}\\[4pt]
Localisation bridge and restriction&
\texttt{Bridge}, \texttt{Restriction}:
\texttt{invertible-le}, \texttt{le-invertible}, \texttt{rho}\\[4pt]
Section rings, basic opens and limits&
\texttt{Sections}, \texttt{BasicOpen}, \texttt{Limit}:
\texttt{Section}, \texttt{iso}, \texttt{natural}, \texttt{lift}\\[4pt]
Basic separation and gluing&
\texttt{Separation}, \texttt{NatBound}, \texttt{Gluing}:
\texttt{injective}, \texttt{Cover.sheaf}\\[4pt]
Sheaf on all finite covers&
\texttt{Separated}, \texttt{RingedLattice}, \texttt{Sheaf}:
\texttt{amalgam-prop}, \texttt{structure-sheaf}\\[4pt]
Extension theorem&
\texttt{Extension}:
\texttt{extension-map}, \texttt{extend-restrict},
\texttt{extend-basic}, \texttt{Unique.agrees}\\[4pt]
Invertibility opens and local structure&
\texttt{InvOpen}, \texttt{SpecInv}, \texttt{LocallyRinged},
\texttt{SpecLocal}: \texttt{InvMap}, \texttt{spec}\\[4pt]
Ring fibre products&
\texttt{FibreProduct}: \texttt{Lift}, componentwise units\\[4pt]
Doubling and the chart identifications&
\texttt{DoubleLattice}, \texttt{Doubling}, \texttt{Chart}:
\texttt{doubled}, \texttt{iota}, \texttt{to}, \texttt{natural},
\texttt{inv-to}, and the second-chart counterparts\\[4pt]
Maps into a doubling&
\texttt{Fold}, \texttt{IntoDoubling}, \texttt{DoublingUniversal}:
\texttt{equivalence}, \texttt{local-equivalence}\\[4pt]
Affine spectrum and doubled-origin line&
\texttt{Affine}, \texttt{DoubledLine}:
\texttt{spec}, \texttt{doubled-line}\\[4pt]
Global sections of the affine spectrum&
\texttt{Global}: \texttt{global}, \texttt{evaluate}, \texttt{iso}\\[4pt]
The universal property of the affine spectrum&
\texttt{SpecMap}, \texttt{SpecUnique}, \texttt{SpecUniversal}:
\texttt{morphism}, \texttt{is-local}, \texttt{lat-eq},
\texttt{sharp-eq}, \texttt{component-of}\\[4pt]
Restriction to an open&
\texttt{Slice}: \texttt{restricted}, \texttt{inclusion},
\texttt{inclusion-local}\\[4pt]
Isomorphisms and their locality&
\texttt{Iso}: \texttt{IsIso}, \texttt{EqualOpens},
\texttt{Transfer}, \texttt{Local.local}\\[4pt]
Chart isomorphisms onto restrictions&
\texttt{ChartIso}: \texttt{chart1}, \texttt{iso1},
\texttt{local1}, and their second-chart counterparts\\[4pt]
Composition and spectra of ring maps&
\texttt{Compose}, \texttt{SpecHom}:
\texttt{Identity}, \texttt{Then.component}, \texttt{Isomorphism.iso}\\[4pt]
The canonical affine criterion&
\texttt{AffineCriterion}: \texttt{IsAffine},
\texttt{Criterion.affine}\\[4pt]
Principal affine opens&
\texttt{SliceUniversal}, \texttt{SpecBasicOpen}:
\texttt{Back}, \texttt{Restore}, \texttt{cancel},
\texttt{j-iso}, \texttt{k-iso}, \texttt{affine}\\[4pt]
Qcqs schemehood and the doubled-origin example&
\texttt{Qcqs}, \texttt{DoubledLineQcqs}:
\texttt{AffineOpen}, \texttt{AffineCover}, \texttt{IsQcqs},
\texttt{affine1}, \texttt{affine2}, \texttt{qcqs}\\[4pt]
Regularity and localisation at the variable&
\texttt{PolynomialCoefficients} (shared from \texttt{sag/pn}),
\texttt{VariableKernel}: \texttt{regular},
\texttt{Kernel.equal}, \texttt{one-zero}\\[4pt]
Global sections of the doubled-origin line&
\texttt{DoubledLineGlobal}: \texttt{separate},
\texttt{equal-components}, \texttt{diagonal},
\texttt{back}, \texttt{iso}\\[4pt]
Non-affineness over a nonzero base&
\texttt{DoubledLineNotAffine}: \texttt{d-diagonal},
\texttt{Image.diagonal}, \texttt{IfAffine.one-zero},
\texttt{not-affine}\\[4pt]
Equality and sets of morphisms&
\texttt{MapEquality}: \texttt{ext}, \texttt{local-ext};
\texttt{SpecUniversal.equivalence}\\[4pt]
Doubled-line points and base change&
\texttt{DoubledLinePoints}: \texttt{equivalence};
\texttt{DoubledLineBaseChange}: \texttt{natural}\\[4pt]
Idempotents and complementary opens&
\texttt{Idempotents}, \texttt{Interval}\\[4pt]
Zariski lattice of a principal quotient&
\texttt{ZariskiQuotient}; shared
\texttt{PrincipalQuotient}, \texttt{PrincipalCertificates}\\[4pt]
Idempotent point classification and naturality&
\texttt{DoubledLineIdempotents},
\texttt{DoubledLineIdempotentsBaseChange}\\[4pt]
Gluing with a transition and the projective line&
\texttt{LatticePullback}, \texttt{GlueAlong},
\texttt{GlueTransition}, \texttt{ProjectiveLine}:
\texttt{affine1}, \texttt{affine2}, \texttt{qcqs}\\[4pt]
Inversion and the punctured chart&
\texttt{Inversion}, \texttt{PuncturedChart}:
\texttt{sigma}, \texttt{involution}, \texttt{automorphism}\\[4pt]
Global functions on the projective line&
\texttt{InversionOverlap}: \texttt{overlap};
\texttt{ProjectiveLineGlobal}: \texttt{diagonal}, \texttt{back},
\texttt{iso}\\[4pt]
Non-affineness of the projective line&
\texttt{ProjectiveLineNotAffine}: \texttt{not-affine}\\[4pt]
Structure map and its chart restrictions&
\texttt{ProjectiveLineStructure}: \texttt{structure},
\texttt{structure-local}, \texttt{First.compatible},
\texttt{Second.compatible}\\
\end{longtable}
\normalsize

The construction of the structure sheaf is parametric in the chosen
principal localisations. \texttt{Affine} supplies the fraction rings as
one such family. Restriction to a principal ideal is packaged as a locally
ringed lattice, and the chart identifications are isomorphisms onto those
restrictions. The affine criterion applies to any principal localisation
family for $A$; the spectrum of $\Gamma X$ used to define the canonical
map uses the fraction-ring family. The formal affine criterion proves the
substantive direction of Theorem~\ref{thm:affine-criterion}; the other
direction is the defining canonical isomorphism itself.

The universal-property proofs first compare lattice maps pointwise and
section maps through restriction between equal target opens.
\texttt{MapEquality} identifies this concrete equality with equality of
the records, and proves that morphisms, and local morphisms, form sets.
The inverse laws for the spectrum and doubling constructions are
therefore exposed as actual equivalences. The spectrum statement remains
parametric in the principal localisation family and explicitly assumes
that the carrier of $A$ is a set.

The doubled-line point classification is proved for $k$-algebras, with
base change both in the extent formulation and in the idempotent
formulation. The quotient-lattice proof uses universal supports; the
complementary-open proof uses sheaf gluing. No synthetic geometric axiom
is used in either.

For the projective line, the checked results include its construction as
a locally ringed lattice, its two affine chart isomorphisms, qcqs
schemehood, the canonical ring isomorphism with $k$, non-affineness when
$1\ne0$ in $k$, and the explicit local structure map to $\Spec k$ with
both chart compatibilities. The proof of
Lemma~\ref{lem:inversion-overlap} is the coefficient-list induction used
in the formal development.

Not yet formalised at this checkpoint are a general mapping property for
gluing with a transition, the classification of maps into the projective
line by extents and reciprocal sections, and its comparison with
rank-one direct summands of $B^2$. No proof of non-separatedness of the
doubled-origin line or separatedness of the projective line is claimed
here. Category-theoretic packaging and comparison with functorial
schemes are outside the present account.

\begin{thebibliography}{9}
\bibitem{cl03}
Thierry Coquand and Henri Lombardi,
\emph{Hidden constructions in abstract algebra: Krull dimension of
distributive lattices and commutative rings},
in \emph{Commutative Ring Theory and Applications},
Lecture Notes in Pure and Applied Mathematics \textbf{231},
Marcel Dekker, New York, 2003, pp.~477--499.
\href{https://arxiv.org/abs/1801.00097}{Updated version: arXiv:1801.00097}.
\bibitem{cls07}
Thierry Coquand, Henri Lombardi and Peter Schuster,
\emph{The projective spectrum as a distributive lattice},
Cahiers de Topologie et G\'eom\'etrie Diff\'erentielle Cat\'egoriques
\textbf{48} (2007), no.~3, 220--228.
\href{https://www.numdam.org/item/CTGDC_2007__48_3_220_0/}{Numdam}.
\bibitem{cls09}
Thierry Coquand, Henri Lombardi and Peter Schuster,
\emph{Spectral schemes as ringed lattices},
Annals of Mathematics and Artificial Intelligence \textbf{56} (2009),
nos.~3--4, 339--360.
\href{https://doi.org/10.1007/s10472-009-9160-7}{doi:10.1007/s10472-009-9160-7}.
\bibitem{zeuner}
Max Zeuner, \emph{Univalent Foundations of Constructive Algebraic
Geometry}. Source copy: \texttt{SCHEME/zeuner1.pdf}.
\bibitem{zm}
Max Zeuner and Anders M\"ortberg,
\emph{A Univalent Formalization of Constructive Affine Schemes},
LIPIcs TYPES 2022, article 14,
\href{https://doi.org/10.4230/LIPIcs.TYPES.2022.14}{doi:10.4230/LIPIcs.TYPES.2022.14}.
Source copy: \texttt{SCHEME/zeuner2.pdf}.
\bibitem{isaev}
Valery Isaev, \emph{A Simple and Constructive Approach to Schemes}.
Source copy: \texttt{SCHEME/isaev.pdf}.
\end{thebibliography}
\end{document}
